\textbf{Mise en situation :} Vous assistez à un match de football du Noir et Or au terrain du Cégep. Alors que l'équipe locale réussit un touché spectaculaire en fin de partie, la sirène de but géante des partisans retentit avec éclat au bord du terrain pour célébrer la première victoire depuis plusieurs années.

\textbf{Questions :} Selon la fiche technique du matériel des partisans du Noir et Or, la puissance sonore émise à la source par cette sirène de but est de $\PneuPuissance\text{ W}$. Si vous mesurez un niveau d'intensité sonore de $\PneuNiveaudB\text{ dB}$ à l'endroit précis où vous êtes assis dans les gradins...

\begin{enumerate}[label=\alph*)]
    \item À quelle distance de la sirène de but vous trouvez-vous ?
    \item Quelle puissance acoustique exacte votre tympan absorbe-t-il si l'on modélise votre conduit auditif par un détecteur circulaire plat de rayon $\OreilleRayon\text{ m}$ ?
\end{enumerate}

Pour en ajouter une couche et faire vibrer la foule, le DJ du stade monte le son des haut-parleurs et fait jouer l'hymne de victoire du Noir et Or. Durant une pause entre deux coups de sirène (lorsque la sirène est silencieuse), vous mesurez le niveau sonore de la musique seule à $\MusiqueNiveaudB\text{ dB}$.

\begin{enumerate}[label= c)]
    \item Quel est le niveau d'intensité sonore total ressenti (en dB) lorsque la sirène de but retentit à nouveau pendant que la musique joue simultanément dans les haut-parleurs ?
\end{enumerate}
